\SourcePage{272}{275}
\section{Partial polarization of particles}

By a suitable choice of direction for the $z$ axis, one component (say $\psi^2$) of any given spinor $\psi^\lambda$, the wave function of a spin-$1/2$ particle, can always be made zero. This is clear already because a spatial direction is specified by two quantities (angles), so the number of parameters at our disposal is precisely the number of quantities we wish to set equal to zero (the real and imaginary parts of the complex number $\psi^2$).

Physically, this means that if a particle of spin $1/2$---for definiteness, let it be an electron---is in a state described by some spin wave function, a direction exists along which the particle's spin projection has the definite value $\sigma=1/2$. The electron in such a state may be said to be \emph{fully polarized}.

There are, however, electron states that may be termed
\emph{partially polarized}. Such states are specified not by wave
functions, but solely by density matrices; in other words, they are mixed
states with respect to spin (see \S~14).

The electron's spin (or \emph{polarization}) density matrix is a rank-two
spinor $\rho^\lambda{}_\mu$, whose normalization is fixed by the condition

\begin{equation}
 \rho^\lambda{}_\lambda=\rho^1{}_1+\rho^2{}_2=1
 \tag{59.1}
\end{equation}
and satisfying the ``Hermiticity'' condition
\begin{equation}
 (\rho^\lambda{}_\mu)^*=\rho^\mu{}_\lambda.
 \tag{59.2}
\end{equation}
For a pure spin state of an electron---that is, a completely polarized
state---the spinor $\rho^\lambda{}_\mu$ reduces to a product of components of
the wavefunction $\psi^\lambda$:

\begin{equation}
 \rho^\lambda{}_\mu=\psi^\lambda\psi_\mu^*.
 \tag{59.3}
\end{equation}

\SourcePage{273}{276}
The diagonal components of the density matrix give the probabilities for the values $+1/2$ and $-1/2$ of the electron spin projection on the $z$ axis. The mean value of this projection is therefore
\[
 \bar s_z=\frac12(\rho^1{}_1-\rho^2{}_2),
\]
or, using (59.1),
\begin{equation}
 \rho^1{}_1=\frac12+\bar s_z,\qquad
 \rho^2{}_2=\frac12-\bar s_z.
 \tag{59.4}
\end{equation}
In a pure state, the mean values of $s_\pm=s_x\pm is_y$ are calculated as
\[
 \bar s_\pm=\psi_\lambda^*\hat s_\pm\psi^\lambda.
\]
Since, according to (55.6) and (55.7), the operators $\hat s_\pm$ are represented by the matrices
\[
 \hat s_+=\begin{pmatrix}0&1\\0&0\end{pmatrix},\qquad
 \hat s_-=\begin{pmatrix}0&0\\1&0\end{pmatrix},
\]
we find
\[
 \bar s_+=\psi_1^*\psi^2,\qquad \bar s_-=\psi_2^*\psi^1.
\]
Accordingly, in a mixed state,
\begin{equation}
 \rho^1{}_2=\bar s_-,\qquad \rho^2{}_1=\bar s_+.
 \tag{59.5}
\end{equation}
With the Pauli matrices, equations (59.4) and (59.5) can be combined in the form
\begin{equation}
 \rho^\lambda{}_\mu=\frac12
 (\delta^\lambda{}_\mu+2\hat{\boldsymbol\sigma}^\lambda{}_\mu\bar{\mathbf s}).
 \tag{59.6}
\end{equation}

Thus every component of the electron polarization density matrix is expressed in terms of the mean values of the components of its spin vector. In other words, the real vector $\bar{\mathbf s}$ completely determines the polarization properties of a spin-$1/2$ particle. In the limiting case of full polarization, one component of this vector, with a suitable choice of axis directions, is $1/2$, while the other two vanish. In the opposite case of an unpolarized state, all three components vanish. In the general case of arbitrary partial polarization and an arbitrary coordinate system, $0\leqslant p\leqslant1$, where
\[
 p=2(\bar s_x^2+\bar s_y^2+\bar s_z^2)^{1/2}
\]
is a quantity that may be called the electron's \emph{degree of polarization}.

\SourcePage{274}{277}
For a particle of arbitrary spin $s$, the density matrix is a spinor $\rho^{\lambda\mu\ldots}{}_{\rho\sigma\ldots}$ of rank $4s$, symmetric in its first $2s$ indices and in its last $2s$ indices, and subject to
\begin{equation}
 \rho^{\lambda\mu\ldots}{}_{\lambda\mu\ldots}=1,
 \tag{59.7}
\end{equation}
\begin{equation}
 (\rho^{\lambda\mu\ldots}{}_{\rho\sigma\ldots})^*
 =\rho^{\rho\sigma\ldots}{}_{\lambda\mu\ldots}.
 \tag{59.8}
\end{equation}
To count the independent components of the density matrix, note that only
$2s+1$ essentially distinct combinations occur among the possible values of
the indices $\lambda,\mu,\ldots$ (and likewise among
$\rho,\sigma,\ldots$). The components of the spinor
$\rho^{\lambda\mu\ldots}{}_{\rho\sigma\ldots}$ are also subject to the single
relation (59.7). Their number is therefore
$(2s+1)^2-1=4s(s+1)$. These components are complex, but relations (59.8) ensure
that this does not enlarge the total number of independent quantities needed
to characterize a particle's state of partial polarization; that number is
thus $4s(s+1)$\footnote{Specifying these quantities is equivalent to specifying
the mean values of the components of the vector $\mathbf s$, together with
all their powers and products of orders 2, 3, \ldots, $2s$ that cannot be
reduced to still lower orders (see Problem~3 of \S~55).}. By comparison, only
$4s$ quantities describe a completely polarized particle: the wave function
$\psi^{\lambda\mu\ldots}$ has $2s+1$ complex components, constrained by one
normalization condition and containing one overall phase that is immaterial
to the state.


Like every spinor of rank $4s$, the spinor $\rho^{\lambda\mu\ldots}{}_{\rho\sigma\ldots}$ is equivalent to a collection of irreducible tensors of ranks $4s$, $4s-2$, \ldots, 0. In the present case there is just one tensor of each of these ranks, since the symmetry properties of $\rho^{\lambda\mu\ldots}{}_{\rho\sigma\ldots}$ allow each contraction to be carried out in only one way: over any one of the indices $\lambda,\mu,\ldots$ and any one of $\rho,\sigma,\ldots$. Moreover, the scalar (rank-zero tensor) is not present as an independent quantity at all, because condition (59.7) reduces it to unity.
